\usepackage[margin=1cm]{geometry}

\setlength{\columnseprule}{1pt}
\setlength{\parindent}{0pt}
\pagenumbering{gobble}

\usepackage{amsmath}
\usepackage{amsfonts}
\usepackage{amssymb}
\usepackage{amsthm}
\usepackage{thmtools}

\usepackage{unicode-math}
\setmainfont{TeX Gyre Pagella}
\setmathfont{TeX Gyre Pagella Math}

\usepackage{stmaryrd}
\usepackage{mathrsfs}
\usepackage{wasysym}

\usepackage[french]{babel}

\newtheoremstyle{myplain}%
  {\topsep}% Space above
  {\topsep}% Space below
  {\itshape}% Body font
  {}% Indent amount
  {}% Head font
  {}% Punctuation after head
  {5pt plus 1pt minus 1pt}% Space after head
  {\bfseries\thmname{#1}\thmnumber{ #2}.{\normalfont\thmnote{ (#3)}}} % Head spec

\newtheoremstyle{mydefinition}%
  {0.4\topsep}% Space above
  {0.4\topsep}% Space below
  {\normalfont}% Body font
  {}% Indent amount
  {}% Head font
  {}% Punctuation after head
  {5pt plus 1pt minus 1pt}% Space after head
  {\bfseries\thmname{#1}\thmnumber{ #2}.{\normalfont\thmnote{ (#3)}}} % Head spec


\declaretheorem[name=Théorème, style=mydefinition]{theo}
\declaretheorem[name=Définition, numberlike=theo, style=mydefinition]{defi}
\declaretheorem[name=Proposition, numberlike=theo, style=mydefinition]{prop}
\declaretheorem[name=Propriété, numberlike=theo, style=mydefinition]{prpr}
\declaretheorem[name=Conjecture, numberlike=theo, style=mydefinition]{conj}
\declaretheorem[name=Corollaire, numberlike=theo, style=mydefinition]{coro}
\declaretheorem[name=Lemme, numberlike=theo, style=mydefinition]{lemm}
\declaretheorem[name=Application, numberlike=theo, style=mydefinition]{appl}

\newtheoremstyle{myremark}%
  {0.5\topsep}% Space above
  {0.5\topsep}% Space below
  {\normalfont}% Body font
  {}% Indent amount
  {}% Head font
  {}% Punctuation after head
  {5pt plus 1pt minus 1pt}% Space after head
  {\itshape\thmname{#1}\thmnumber{ #2}.{\normalfont\thmnote{ (#3)}}} % Head spec
\declaretheorem[style=mydefinition,name=Remarque, numberlike=theo]{rema}
\declaretheorem[style=mydefinition,name=Exemple, numberlike=theo]{exem}
\declaretheorem[style=mydefinition,name=Convention, numberlike=theo]{conv}

\newcommand{\R}{\mathbb{R}}
\newcommand{\Rd}{\mathbb{R}^d}
\newcommand{\N}{\mathbb{N}}
\newcommand{\Z}{\mathbb{Z}}
\newcommand{\Q}{\mathbb{Q}}
\newcommand{\K}{\mathbb{K}}
\newcommand{\C}{\mathbb{C}}
\newcommand{\Ps}{\mathfrak{P}}
\newcommand{\dd}{\mathrm{d}}
\newcommand{\Tn}{\mathbb{T}^n}
\DeclareMathOperator{\imm}{im}
\DeclareMathOperator{\supp}{supp}
\DeclareMathOperator{\diam}{diam}
\DeclareMathOperator{\sgn}{sgn}
\DeclareMathOperator{\Aut}{Aut}
\DeclareMathOperator{\Gal}{Gal}
\DeclareMathOperator{\hess}{Hess}
\DeclareMathOperator{\Id}{Id}
\newcommand{\llb}{[\!\![}
\newcommand{\rrb}{]\!\!]}
\newcommand{\ol}[1]{\overline{#1}}
\newcommand{\ub}[2]{\underbrace{#1}_{#2}}
\newcommand{\wt}[1]{\widetilde{#1}}
\newcommand{\wh}[1]{\widehat{#1}}
\newcommand{\eqdef}{\stackrel{\mathrm{def}}{=}}
\newcommand{\wout}[1]{\setminus \{#1\}}
\newcommand{\abs}[1]{\left\lvert #1 \right\rvert}
\newcommand{\norm}[1]{\left\lVert #1 \right\rVert}
\newcommand{\normsub}[1]{\left\lVert \left\lvert #1 \right\rvert \right\rVert}

\AtBeginDocument{%à cause d'unicode-math
\let\leq\leqslant
\let\geq\geqslant
\let\epsilon\varepsilon
}

\usepackage[style=alphabetic]{biblatex}
\bibliography{../my}